\section{Introduction}\label{sec:short_intro}

Option prices depend on parameters that are estimated rather than observed. A Bayesian analysis can account for this uncertainty by integrating the conditional option value over the posterior distribution instead of fixing all parameters at point estimates. This idea is established in Bayesian option pricing and forecasting \cite{guidolin2003,romboutsstentoft2014,guptareisinger2014}. The practical question is narrower: when does that integration materially change the option value?

Arithmetic-average Asian options provide a useful setting because their exposure to volatility changes with moneyness, maturity, and the fraction of the averaging period already fixed. The numerical literature has developed efficient methods for this path dependence, including geometric-average control variates and conditioning approximations \cite{kemnavorst1990,curran1994}. Commodity applications additionally require a term structure of futures prices and, in richer specifications, stochastic convenience yields, jumps, or stochastic volatility \cite{schwartz1997,shirayatakahashi2011,ewaldwu2023}. The objective here is not to introduce another Asian-option model. It is to isolate the effect of parameter uncertainty within a transparent pricing model and to compare that effect with the empirical importance of the volatility information supplied to the model.

The distinction matters because two different questions are easily conflated. The first is whether averaging the conditional option value over a posterior distribution changes the posterior mean price. The second is whether the volatility input itself is aligned with the volatility information embedded in option prices. These are conceptually separate. The first is a nonlinear-averaging problem. The second concerns model specification and information content.

Let \(C^Q(\sigma)\) be the risk-neutral option value conditional on volatility \(\sigma\), and let \(\bar\sigma=\mathbb E[\sigma\mid\mathcal D]\). The posterior-integrated and posterior-mean prices are
\begin{equation}
 C_{PI}=\mathbb E[C^Q(\sigma)\mid\mathcal D],
 \qquad
 C_{PM}=C^Q(\bar\sigma).
 \label{eq:short_pi_pm}
\end{equation}
A second-order expansion gives
\begin{equation}
 C_{PI}-C_{PM}
 \approx
 \frac12 C^{Q\prime\prime}(\bar\sigma)
 \operatorname{Var}(\sigma\mid\mathcal D).
 \label{eq:short_taylor}
\end{equation}
Posterior integration therefore matters when parameter uncertainty coincides with curvature of the option-pricing function. By contrast, the dispersion of the posterior distribution of model prices is locally first order in the posterior standard deviation. A small PI--PM difference need not imply little uncertainty in the theoretical option value.

The empirical application uses WTI Average Price Options (APOs), whose payoff depends on the arithmetic average of daily first-nearby WTI futures settlements. WTI is useful for this comparison because historical volatility can change sharply, the futures curve varies across valuation dates, and option-implied volatility is observable from both the APO market and a separate vanilla-WTI option market. The design uses only information available before each target date when evaluating out-of-sample pricing performance.

Three findings organize the paper. First, a large synthetic experiment confirms Equation~\eqref{eq:short_taylor}: posterior integration becomes important primarily in weak-information states with substantial volatility curvature and little of the averaging period already fixed. Second, the observed October-2026 WTI sample lies in a region where posterior-integrated and posterior-mean prices are almost identical, even though the posterior distribution of model prices remains materially dispersed. Third, prior option-implied volatility inputs produce substantially lower out-of-sample pricing errors than the historical-volatility benchmarks evaluated here. Rolling historical windows, EWMA, and a horizon-matched GARCH(1,1) forecast do not reproduce that performance. A same-contract lagged-IV benchmark shows that the result does not require a flexible cross-sectional smile, and a separate validation using vanilla-WTI options shows that it is not confined to volatility inferred from the APO family itself.

These results do not identify a pure structural difference between physical and risk-neutral volatility. Implied volatility under a parsimonious pricing model can absorb risk premia, stochastic volatility, jumps, omitted futures-curve dynamics, liquidity, and settlement conventions. The supported statement is therefore empirical and conditional: in this sample, prior option-implied volatility contains substantially more information for subsequent APO settlement pricing than the historical-volatility inputs considered here, while posterior integration of the historical-volatility distribution changes the point-price mean much less.

The remainder of the paper is organized around those three findings. Section~\ref{sec:short_method} develops the posterior-pricing mechanism. Section~\ref{sec:short_design} describes the WTI data and out-of-sample design. Section~\ref{sec:short_results} reports the synthetic and empirical results. Section~\ref{sec:short_robustness} summarizes robustness and interpretation, with implementation details and extended diagnostics reported in the online supplement. Section~\ref{sec:short_conclusion} concludes.
